%%%PAGE 43 rot=0 legibility=good summary=几何关联(V是三角形直角边)：楔块顶起重物、球与斜面体、A楔块与B弹簧三道关联速度+能量题；红字“受力要整体法”；末题[4]弹簧+摩擦的两物体加速度与分离条件
%%%BLOCK id=043-01 ch=04 sec=关联速度·几何关联 kind=knowledge alt=none cont=none y=0.02-0.07
\begin{keybox}[几何关联]% 红笔框(页面左上角)，框内两行
$V$ 是三角形直角边
\end{keybox}
%%%END
%%%BLOCK id=043-02 ch=04 sec=关联速度·几何关联 kind=note alt=02,03 cont=none y=0.02-0.06
\red{受力要整体法}% 红色大字，位于页面右上角
%%%END
%%%BLOCK id=043-03 ch=04 sec=关联速度·几何关联 kind=problem alt=02,06 cont=none y=0.02-0.19
\begin{problem}[{[1]}]
%FIG p043_a 0.04 0.015 0.29 0.13
\notefig[0.3]{figures/p043_a.png}
% 图内标注：图左上角红框“几何关联 / V是三角形直角边”(已在上一块转录)；楔形斜面(左下角 $\theta$)，斜面上顶着一根竖直杆，杆顶标 $M$(重物)；楔形右侧标 $m$，水平向左箭头 $F$；红字 $V_M$(竖直向下，右侧)、$V_m$(水平，楔块下方红线箭头)
① 为将重物顶起，求 $F_{\min}$

\noindent $F_N=\dfrac{Mg}{\cos\theta}$\qquad $F_N\sin\theta=\dfrac{Mg}{\cos\theta}\sin\theta$

\noindent 将小轮推高地面 $h$ 处，突然撤去 $F$，求当小轮到水平面时斜面速度

\noindent ②
\[
\left\{\begin{aligned}
&V_M=V_m\tan\theta\\
&Mgh=\tfrac12MV_M^{2}+\tfrac12mV_m^{2}
\end{aligned}\right.
\]
\end{problem}
%%%END
%%%BLOCK id=043-04 ch=04 sec=关联速度·几何关联 kind=problem alt=02,06 cont=none y=0.19-0.44
\begin{problem}[{[2]}]
%FIG p043_b 0.09 0.19 0.345 0.325
\notefig[0.3]{figures/p043_b.png}
% 图内标注：竖直墙(左，顶端标 1)与斜面体(倾角 $30^\circ$，斜面上方标 2、$u=0$)之间夹一个小球(半径 $R$，标 $m$、两处 $R$)，球心到切点的虚线；斜面体右侧竖直边标 $M$，向左箭头 $F$，右侧红色小箭头 $V_1$，底边红箭头 $V_2$，右下标 $u=0$
撤去 $F$ 后小球和斜面体做匀加速直线运动，直到小球落地。

\noindent ① 求 $F$\quad ② 球落地时求 $V_1$、$V_2$

\noindent $F=N=\dfrac{\sqrt3}{3}mg$

\noindent $mgh=\tfrac12mV_1^{2}+\tfrac12mV_2^{2}$

\noindent $h_0=\sqrt3R$\quad $h=h_0-R=(\sqrt3-1)R$

\noindent $x_1=x_2\tan\theta$

\noindent $V_2=\sqrt3V_1$

\noindent $V_1=\sqrt{\dfrac{(\sqrt3-1)gR}{2}}$\quad $V_2=\sqrt{\dfrac{3(\sqrt3-1)gR}{2}}$
\end{problem}
%%%END
%%%BLOCK id=043-05 ch=04 sec=关联速度·几何关联 kind=problem alt=02,06 cont=none y=0.44-0.78
\begin{problem}[{[3]}]
%FIG p043_c 0.065 0.447 0.335 0.62
\notefig[0.3]{figures/p043_c.png}
% 图内标注：左侧楔形物块 $A$(顶角处标 $m$，下角 $37^\circ$、$\alpha$，红字 $V_B$ 水平向右箭头在顶端，红字 $V_A$ 竖直向下箭头在左侧)，紧靠的方块 $B$(标 $2m$)，弹簧(标 $k$)连接 $B$ 与右侧竖直墙；图下标 $\mu=0$；图左上角有题号 [3]
$B$ 将弹簧压缩 $x_0$ 后 $A$ 上滑，弹簧恢复原长时，$B$ 速度为 $V_B$。

\noindent ① 求刚恢复原长时 $V_A$。

\noindent $V_B=V_A\tan\alpha$\qquad $V_A=\dfrac43V_B$

%FIG p043_d 0.635 0.50 0.715 0.562
\notefig[0.25]{figures/p043_d.png}
% 图内标注：矢量图——向上箭头 $V_A$，向左箭头 $V_B$，夹角 $37^\circ$，虚线(斜向左上)
\noindent ② 求弹簧压缩量为 $x_0$ 时具有的弹性势能

\noindent 弹簧原长时，$h_A=\dfrac{x_0}{\tan\alpha}=\dfrac43x_0$

\noindent $E_p=\tfrac12mV_A^{2}+\dfrac{2}{2}mV_B^{2}+mgh=\dfrac{17}{9}mV_B^{2}+\dfrac43mgx_0$

\noindent ③ 求两物体动能最大时，弹簧压缩量

\noindent 设$AB$相互作用力$F$
\[
\left\{\begin{aligned}
&a_A=a_B=0\\
&F\sin\alpha=mg\\
&F\cos\alpha=kx
\end{aligned}\right.\qquad x=\frac{mg}{k\tan\alpha}=\frac{4mg}{3k}
\]
\end{problem}
%%%END
%%%BLOCK id=043-06 ch=03 sec=连接体·弹簧与分离 kind=problem alt=06 cont=none y=0.78-1.00
\begin{problem}[{[4]}]
%FIG p043_e 0.065 0.78 0.31 0.862
\notefig[0.3]{figures/p043_e.png}
% 图内标注：水平向右箭头 $F$ 推方块 $B$，$B$ 紧靠方块 $A$(两者上方各标 $m$)，$A$ 右侧连弹簧，弹簧另一端接右侧竖直墙；地面有斜线(粗糙)，弹簧下方标 $\mu$；图左上角有题号 [4]
$F$ 缓慢推动 $B$，压缩弹簧 $x_0$。此时 $A$、$B$ 静止，$B$ 不受摩擦。撤去 $F$ 后，$A$ 运动最大距离为 $4x_0$。求撤去 $F$ 瞬间，$A$、$B$ 的加速度。

\noindent $kx_0-\mu mg=2ma$% 原稿此处 k 写作大写 K

\noindent $a=\dfrac{kx_0}{2m}-\dfrac{\mu g}{2}$

\noindent $A$、$B$ 一起向左运动 $x$ 后\underline{分离}，求 $x$% “分离”二字下有红色下划线

\noindent 分离临界 $N=0$\quad $a_A=a_B=0$\qquad $k(x_0-x)=\mu mg$\qquad $\therefore x=x_0-\dfrac{\mu mg}{k}$
\end{problem}
%%%END
%%%PAGEEND 43
